% ======================================================================
\section{Collinear (DGLAP) and running-coupling terms}\label{sec:dglap_rc}
% ======================================================================
% Sources: paper/notes/fix_poles_dglap.tex (P3, P4, P5, Row II.IV tail), notes/lit_dglap_rc.tex,
% limits_checked/Finite_terms_final.tex (Group II Row IV), limits_checked/Divergent_terms_checked.tex.

In this section we discuss the terms that arise when the integrated gluon of Row~II.IV is collinear with the observed one, and we give the account of all poles $1/\epsilon$ of the calculation. Although, following the convention of the rest of the paper, no pole is written in the final result, the poles determine what the collinear and the running-coupling terms are. We use the notation
\begin{align}\label{eq:dg_not}
\mathcal O(y,y')&\equiv\Kv(x-y)\cdot\Kv(x'-y')\,\big[U(y')-U(x')\big]^{ce'}\big[U(y)-U(x)\big]^{ce}\rho^{e'}(x')\,\rho^{e}(x)\,,\notag\\
\widehat L(q)&\equiv\int_{x,y,x',y'}e^{-iq\cdot(y-y')}\,\mathcal O(y,y')\,,\notag\\
\frac{d^3N_{\rm LO}}{d^3k}&=\mathcal N_0\,\widehat L(k)\,,\qquad
\mathcal N_0\equiv\frac{1}{(2\pi)^3}\frac{g^2}{2\pi^2}\frac{1}{k^+}\,,\notag\\
\mathcal C&\equiv\frac{\alpha_s}{2\pi}\,\mathcal N_0=\frac{1}{(2\pi)^3}\frac{g^4}{16\pi^4}\frac{1}{k^+}\,,
\end{align}
so that Eq.~\eqref{eq:LO} is $\mathcal N_0\widehat L(k)$, and
\begin{align}\label{eq:dg_B}
\mathcal B(r)\equiv2\gamma_E+\log\frac{r^2\mu^2}{4}=\log\frac{r^2\mu^2}{c_0^2}\,,\qquad c_0=2e^{-\gamma_E}\,,\qquad
\mathfrak a\equiv\frac{\alpha_sN_c}{\pi}\,,\qquad \frac{\alpha_s\beta_0}{4\pi}=\frac{11}{12}\,\mathfrak a\,,
\end{align}
with $\mu\equiv\mu_{\MS}$, $\beta_0=11N_c/3$, $L\equiv\log(k^+/\Lambda)$ and $\ell_\vee\equiv\log(1-k^+/\vee)$. The phase is $e^{-i\v k\cdot(\text{unprimed}-\text{primed})}$ throughout (Section~\ref{sec:fin_conv}); in the collinear terms it is $e^{-i\v k\cdot(\v y-\v y')/\xi}$.

\subsection{Origin: final-state splitting of the LO gluon}\label{sec:dg_origin}

Row~II.IV is the square of the amplitude $\mathcal A_{CA}$ of Eq.~\eqref{eq:A2gamp}: a gluon of momentum $p^++k^+$ produced at LO splits after the collision, through $\bar{\mf C}_1$, into the observed gluon (at $\v w$ in the amplitude, $\v w'$ in the conjugate amplitude) and a second gluon at $\v z$, which is integrated over. With $\xi\equiv k^+/(p^++k^+)$ the momentum fraction of the observed gluon in its parent, the parent sits at $\v y=\v z+\xi(\v w-\v z)$ and $\v y'=\v z+\xi(\v w'-\v z)$; the Wilson line $U(\v z)$ cancels by unitarity, and the parent carries the Wilson-line structure of the LO amplitude. The integral over the daughter separation $\v r=\v w-\v z$ at fixed $\v u=\v w'-\v w$ diverges only at large $|\v r|$, i.e.\ at small relative transverse momentum of the two daughters: it is a collinear (infrared) divergence. Its coefficient is the real splitting function
\begin{align}\label{eq:Pggplus}
\mathcal P_{gg+}(\xi)=2N_c\Big[\xi(1-\xi)+\frac{\xi}{1-\xi}+\frac{1-\xi}{\xi}\Big]\,,
\end{align}
without a $\delta(1-\xi)$ term, as appropriate for a real emission. The two end points are the rapidity logarithms of the row: $\xi\to1$ (the integrated gluon soft) gives its JIMWLK term, $\xi\to0$ (the observed gluon soft with respect to its parent) its BFKL term (Section~\ref{sec:rapidity}). It is convenient to write the LO structure of a parent gluon of transverse momentum $\v k/\xi$ as $\Pi(\xi)\equiv\xi^{-2}\widehat L(\v k/\xi)$, and to define
\begin{align}\label{eq:dg_X}
\mathcal X\equiv\int_{x,x',z}\mathcal O(z,z)=\int_{x,x',z}&\Kv(x-z)\cdot\Kv(x'-z)\notag\\&\times\big[U(z)-U(x')\big]^{ce'}\big[U(z)-U(x)\big]^{ce}\rho^{e'}(x')\,\rho^e(x)\,,
\end{align}
the structure that multiplies the BFKL logarithm of this row, $+2N_c\,\mathcal C\,\frac{4\pi}{\v k^2}\log\frac{\vee}{k^+}\,\mathcal X$.

\subsection{Account of the poles}\label{sec:dg_poles}

We regularise the transverse loop and collinear integrals in $d=2-\epsilon$ dimensions. A scaleless integral is $\int d^dp/\v p^2=\pi\big(\frac{2}{\epsilon_{\rm UV}}-\frac{2}{\epsilon_{\rm IR}}\big)$; with a single $\epsilon$, $\epsilon_{\rm UV}=\epsilon_{\rm IR}$, it vanishes, which is the convention used elsewhere in the paper. We give both bookkeepings; they lead to the same finite result. Four sources of poles have to be considered.

\paragraph{Scheme.} We use conventional dimensional regularisation (CDR): transverse momenta and gluon polarisations are continued to $d_\perp=2-\epsilon$ dimensions wherever a polarisation sum multiplies a pole, i.e.\ in Rows~I.II, I.V and II.IV; in the finite terms, which are locally integrable, $d_\perp=2$. The coupling is the $\MS$ coupling and the fragmentation function the $\MS$ number density (pure-pole counterterms with the four-dimensional $\mathcal P_{gg}$, Section~\ref{sec:dg_msbar}), both in this scheme; Eq.~\eqref{eq:A3onerho} and the constant $\frac{67}{9}-\frac{\pi^2}{3}$ refer to it. In the dilute limit the two-gluon amplitude of the real rows reproduces, for any number $d_\perp$ of transverse polarisations, the $(D-2)$ dependence of the CDR two-gluon production cross section of~\cite{Fadin:1998py,Bartels:2006hg} (Section~\ref{sec:check_dilute}). With two-dimensional polarisations and $d$-dimensional momenta (a dimensional-reduction-like scheme) the reduction $\tilde p^l\tilde p^i\to\delta^{li}\tilde p^2/d$ gives a factor $2/d$ instead of $d_\perp/d=1$: the constant of Row~I.II would be $\frac{64}{9}-\frac{\pi^2}{3}$ (a shift $-\frac{\alpha_sN_c}{6\pi}\times$LO), that of Row~I.V would change by the opposite amount, as it must by Eq.~\eqref{eq:dg_SEid} below, Row~II.IV by $-\mathcal C\int_0^1d\xi\,2N_c\,\xi(1-\xi)\,\Pi(\xi)$, and the coupling and the fragmentation function would have to be converted, $\alpha_s^{\rm DR}=\alpha_s^{\MS}\big(1+\frac{\alpha_sN_c}{12\pi}\big)$ and $D^{\rm DR}=D^{\MS}\otimes\big[\delta(1-\xi)+\frac{\alpha_s}{2\pi}\big(2N_c\xi(1-\xi)-\frac{N_c}{6}\delta(1-\xi)\big)\big]$; the coefficient $\frac{N_c}{6}$ of the $\delta$ term is fixed independently by the momentum sum rule and equals the coupling conversion. With these conversions the result is unchanged; the shift of order $N_c/6$ cancels in the virtual part and appears only in the conversion of $\alpha_s$ and $D$. (The sign of the Row~II.IV term was obtained from the general $d_\perp/d$ rule and from this closure, not by re-deriving the row.)

\paragraph{The vertex correction (Row I.II).} The one-$\rho$ part of $\mf A^{(3)}$, Eq.~\eqref{eq:A3onerho}, carries $\big(\frac{11}{3}-2L+\ell_\vee\big)\big[-\frac{2}{\epsilon}+\log\frac{\v k^2}{\mu^2}\big]$ times $\mf A^{(1)}$; in coordinate space $\log\frac{\v k^2}{\mu^2}\frac{k^i}{\v k^2}\to-\mathcal B(\v x-\v w)\frac{(x-w)^i}{(\v x-\v w)^2}$, Eq.~\eqref{eq:FT_log}. In the cross section, with its complex conjugate, the $\frac{11}{3}$ part gives
\begin{align}\label{eq:dg_II}
\frac{d^3N^{\rm I.II}_{11/3}}{d^3k}=\frac{\alpha_s\beta_0}{4\pi}\,\mathcal N_0\int_{x,w,x',w'}e^{-ik\cdot(w-w')}\Big[\frac{4}{\epsilon}+\mathcal B(x-w)+\mathcal B(x'-w')\Big]\,\mathcal O(w,w')\,.
\end{align}
This is the virtual term $\frac{\alpha_s\beta_0}{4\pi}\big[\frac4\epsilon+L(\v x-\v y)+L(\v x'-\v y')\big]\times$LO expected from the self-energy of the emitted gluon. The pole comes from the region of large loop momentum of the $gg$ loop, $\int d^2\tilde p\,\tilde p^2/(\tilde p^2+\xi(1-\xi)\v k^2)$, which is finite at small $\tilde p$: it is ultraviolet.

\paragraph{The recombination term (Row I.V).} The part of $\bar{\mf C}_1^\dagger\mf B_2$ in which the two gluons of $\mf B_2$ are at one point (``$N_{2b\,\rm UV}$'') has the longitudinal coefficient $\frac{1}{k^+}\int_\Lambda^{k^+-\Lambda}dp^+[\dots]=2L-\frac{11}{6}$ and the bracket $\frac2\epsilon+\mathcal B(\v y-\v x)$, the pole coming from $\v z\to\v x$, again ultraviolet (the $\v z$ integral has no large-distance divergence). After the relabelling $\v x\to\v w$, $\v y\to\v x$, which turns the phase into $e^{-i\v k\cdot(\v w-\v w')}$ and leaves the order of the charges unchanged, and with $\frac{g^4N_c}{16\pi^4}\frac{11}{6}=\frac{g^2}{2\pi^2}\frac{\alpha_s\beta_0}{4\pi}$, its $-\frac{11}{6}$ part is
\begin{align}\label{eq:dg_V}
\frac{d^3N^{\rm I.V}_{2b\,\rm UV}}{d^3k}\bigg|_{-\frac{11}{6}}&=-\frac{\alpha_s\beta_0}{4\pi}\,\mathcal N_0\int_{x,w,x',w'}e^{-ik\cdot(w-w')}\Big[\frac{4}{\epsilon}+\mathcal B(x-w)+\mathcal B(x'-w')\Big]\,\mathcal O(w,w')\notag\\&=-\frac{d^3N^{\rm I.II}_{11/3}}{d^3k}\,,
\end{align}
pole, $2\gamma_E$ and logarithm, provided the same $\mu$ is used. The $-\frac{11}{6}$ of $N_{2b\,\rm UV}$ is the dimensionally regularised $\v z$ integral of a subtraction term; the same term, unintegrated and with the opposite sign, remains in the finite part of the row (in $N_{2b\,\rm rest}$), where it is the term with $\frac{11}{6}$ of Eq.~\eqref{eq:S-I.V-2o}, and makes the point $\v z\to\v x$ integrable; its $\v z$ integral alone diverges. The cancellation against Row~I.II concerns the integrated subtraction only. The cancellation of Eqs.~\eqref{eq:dg_II} and~\eqref{eq:dg_V} is a cancellation of an ultraviolet pole against an ultraviolet pole. Its physical origin is that a point-like $gg$ fluctuation of the produced gluon does not resolve the shock wave, $f^{adc}f^{fbe}U^{db}(\v x)U^{ce}(\v z)\to N_cU^{af}(\v x)$ at $\v z=\v x$. The observable counts dressed gluons, $\Omega^\dagger S^\dagger\Omega\,a^\dagger a\,\Omega^\dagger S\Omega$, and where $S$ is locally constant $\Omega^\dagger S\Omega=S$: the ultraviolet part of the self-energy before the shock wave (Row~I.II) is undone by the same loop straddling it (Row~I.V). 

\paragraph{Result and grouping.} The cancellation of the poles and of the dependence on $\mu$ between Rows~I.II and~I.V is a result, for any grouping of the finite terms. Which finite terms accompany it depends on the grouping. Let $a$ be the light-cone energy of the emitted gluon, $D$ the relative energy of the gluon pair into which it fluctuates and $b=a+D$ the energy of the pair. The self-energy part of $\mf A^{(3)}$ in Row~I.II has the energy denominators $1/(ab)$; the part of Row~I.V in which the pair produced by the triple-gluon vertex (together with the instantaneous term, which cancels the longitudinal part of the polarisation sum) is recombined by $\bar{\mf C}_1^\dagger$ without resolving the target has $1/(bD)$, with the same vertices and the same longitudinal supports $\Lambda<p^+<k^+-\Lambda$; and
\begin{align}\label{eq:dg_SEid}
\frac{1}{ab}+\frac{1}{bD}=\frac{1}{aD}\,,
\end{align}
whose transverse integral is scaleless and vanishes in dimensional regularisation, for each momentum fraction $\xi$, in $d$ dimensions. Hence, if the colour of the pair in Row~I.V is evaluated at its transverse centre of mass $\xi\v x+(1-\xi)\v z$, the self-energy part of Row~I.II, with its pole, its logarithm $\mathcal B$, its soft end points and its constant $\frac{67}{9}-\frac{\pi^2}{3}$, is cancelled identically by the unresolved part of Row~I.V, and the $\beta_0$-type coefficient survives only in the region in which the pair resolves the target. (The commutator parts of $\mf A^{(3)}$ and of $\mf B_2$ are not covered by Eq.~\eqref{eq:dg_SEid}.) In the form used in this paper the colour of the pair is evaluated at the position $\v x$ of one of its gluons; then no term $\frac{\alpha_s\beta_0}{4\pi}[\mathcal B(\v x-\v w)+\mathcal B(\v x'-\v w')]\times$LO remains explicitly, and the region in which the pair resolves the target is in the finite terms of Row~I.V ($N_{2a}$, $N_{2b\,\rm rest}$), which carry the same coefficient $2L-\frac{11}{6}$ multiplying the small-pair kernel. For an emitter distance $|\v x-\v w|$ larger than the inverse saturation scale of the target they contain, by power counting of the kernel, a logarithm of $|\v x-\v w|Q_s$ with this coefficient (not evaluated with a target model). The statement that the emission coupling does not run with the distance $\v x-\v w$ between the emitter and the produced gluon is therefore a statement about this grouping, valid as long as the pair does not resolve the target within this distance. In NLO JIMWLK, where the soft gluon is not measured, the $\beta_0$ terms with the emitter distance survive~\cite{Kovner:2023vsy}. The identity~\eqref{eq:dg_SEid} was confirmed independently in momentum space, in the dilute limit with a point-like projectile, where the unresolved part of Row~I.V cancels the whole self-energy part of Row~I.II $\xi$ by $\xi$, to $10^{-13}$ (Section~\ref{sec:check_dilute}).

\paragraph{The normalisation of the dressed observed gluon.} The single-gluon normalisation contains the scaleless integral $\braket{\Psi_{gg}|\Psi_{gg}}=\frac{g^2N_c}{8\pi^3}\big(2L-\frac{11}{6}\big)\int\frac{d^2\tilde p}{\tilde p^2}$, which is set to zero by the convention above. In the cross section it enters as
\begin{align}\label{eq:dg_norm}
\frac{d^3N_{\rm norm}}{d^3k}=-\frac{g^2N_c}{8\pi^3}\Big(2L-\frac{11}{6}\Big)\int\frac{d^d\tilde p}{\tilde p^2}\;\frac{d^3N_{\rm LO}}{d^3k}
=\Big[\frac{\alpha_s\beta_0}{4\pi}-\mathfrak aL\Big]\Big(\frac{2}{\epsilon_{\rm UV}}-\frac{2}{\epsilon_{\rm IR}}\Big)\frac{d^3N_{\rm LO}}{d^3k}\,.
\end{align}
Its ultraviolet half, $\frac{\alpha_sN_c}{4\pi}\big(\frac{11}{3}-4L\big)\frac{2}{\epsilon_{\rm UV}}$, is the light-cone-gauge $Z_3-1$ of the produced gluon; its infrared half ($\tilde p\to0$, collinear pair) is the virtual partner of the real splitting of Row~II.IV.

\paragraph{The real splitting (Row II.IV).} After the extraction of the end points, the collinear poles of Row~II.IV are: in its JIMWLK term $-(L+\ell_\vee)\,\mathfrak a\,\frac2\epsilon\,\frac{d^3N_{\rm LO}}{d^3k}$; in its DGLAP term
\begin{align}\label{eq:dg_IVpole}
-\mathcal C\,\frac2\epsilon\int_0^1d\xi\,\mathcal P_{gg+}(\xi)\,\Pi(\xi)=-\mathcal C\,\frac2\epsilon\,\big[\mathcal P_{gg}\otimes\Pi\big]+\frac{\alpha_s\beta_0}{4\pi}\,\frac2\epsilon\,\frac{d^3N_{\rm LO}}{d^3k}\,,
\end{align}
where $\mathcal P_{gg}=\mathcal P_{gg+}+\frac{\beta_0}{2}\delta(1-\xi)$ is the full $\MS$ splitting function at $N_f=0$, with the plus prescriptions of Eq.~\eqref{eq:plusxi}, and $\mathcal C\frac{\beta_0}{2}\Pi(1)=\frac{\alpha_s\beta_0}{4\pi}\mathcal N_0\widehat L(k)$. The subtraction at $\xi\to0$ has no pole, since the coefficient of the pole vanishes at $\xi=0$ for $\v k\neq0$. The second term of Eq.~\eqref{eq:dg_IVpole} is the pole of what we call the running-coupling term below: it is collinear, not ultraviolet.

\paragraph{Counterterms.} The coupling is renormalised in the $\MS$ scheme, $g_0^2=\mu^\epsilon g^2(\mu)\big[1-\frac{\alpha_s\beta_0}{4\pi}\frac2\epsilon\big]$; in the LO spectrum this gives $-\frac{\alpha_s\beta_0}{4\pi}\frac{2}{\epsilon_{\rm UV}}\times$LO. The charges $\rho$ and the Wilson lines $U$ are not renormalised: the source is classical and $U$ is the given target field. The collinear pole of the observed gluon is absorbed into an $\MS$ fragmentation function with the full $\mathcal P_{gg}$, which gives the counterterm $+\mathcal C\frac{2}{\epsilon_{\rm IR}}[\mathcal P_{gg}\otimes\Pi]$, with the $d$-dimensional Jacobian of Section~\ref{sec:dg_msbar}.

\begin{table}[t]
\centering
\renewcommand{\arraystretch}{1.25}
\footnotesize
\begin{tabular}{@{}p{3.1cm}p{2.6cm}p{0.6cm}p{2.4cm}p{4.0cm}@{}}
\toprule
pole & region & type & coefficient & cancelled by / absorbed into\\
\midrule
Row I.II, $\mf A^{(3)}$ & $gg$ loop, $\tilde p\to\infty$ & UV & $\frac{11}{6}-L+\frac{\ell_\vee}{2}$ & $\frac{11}{6}$: Row I.V; $-L$: JIMWLK sector\\
Row I.V, $N_{2b\,\rm UV}$ & pair at one point, $\v z\to\v x$ & UV & $-\frac{11}{6}+2L$ & $-\frac{11}{6}$: Row I.II; $2L$: JIMWLK sector\\
normalisation, UV half & $\tilde p\to\infty$ & UV & $\frac{11}{12}-L$ & $\frac{11}{12}$: coupling counterterm; $-L$: JIMWLK sector\\
coupling counterterm & $g_0\to g(\mu)$ & UV & $-\frac{11}{12}$ & UV half of the normalisation\\
\midrule
normalisation, IR half & collinear pair, $\tilde p\to0$ & IR & $-\frac{11}{12}+L$ & $-\frac{11}{12}$: $\delta(1-\xi)$ pole of Row II.IV; $L$: Row II.IV, $\xi\to1$\\
Row II.IV, $\xi\to1$ (JIMWLK) & large $\v r$ & IR & $-L-\ell_\vee$ & $-L$: normalisation, IR half; $-\ell_\vee$: left\\
Row II.IV, real $\mathcal P_{gg+}$ & large $\v r$ & IR & $-[\mathcal P_{gg}\otimes\Pi]+\frac{11}{12}$ & $-[\mathcal P_{gg}\otimes\Pi]$: FF counterterm; $\frac{11}{12}$: normalisation, IR half\\
Row II.IV, $\xi\to0$ (BFKL) & --- & --- & $0$ & ---\\
FF counterterm, full $\mathcal P_{gg}$ & $D_g\to D_g(\mu)$ & IR & $+[\mathcal P_{gg}\otimes\Pi]$ & real $\mathcal P_{gg+}$ and virtual $\frac{\beta_0}{2}\delta(1-\xi)$\\
\midrule
sum & & & UV: $\frac{\ell_\vee}{2}$;\ IR: $-\ell_\vee$ & $\to0$ for $\vee\to\infty$\\
\bottomrule
\end{tabular}
\caption{The poles of the calculation. Coefficients in units of $\mathfrak a\,\frac2\epsilon\,\frac{d^3N_{\rm LO}}{d^3k}$, and of $\mathcal C\frac2\epsilon[\,\cdot\otimes\Pi]$ for the collinear (DGLAP) entries; complex conjugates included for Group~I; $\mathfrak a=\alpha_sN_c/\pi$, so that $\frac{11}{12}\mathfrak a=\frac{\alpha_s\beta_0}{4\pi}$. The upper block contains the ultraviolet poles, the lower block the infrared (collinear) ones.}
\label{tab:poles}
\end{table}

\paragraph{The balance.} The poles are collected in Table~\ref{tab:poles}; the coefficients were obtained in exact rational arithmetic. In the single-$\epsilon$ bookkeeping, in which the normalisation~\eqref{eq:dg_norm} vanishes, the sum is Row~I.II $+$ Row~I.V $+$ Row~II.IV $+$ coupling counterterm $+$ fragmentation counterterm $=\big(\frac{11}{6}-L+\frac{\ell_\vee}{2}\big)+\big(-\frac{11}{6}+2L\big)+\big(-L-\ell_\vee+\frac{11}{12}-[\mathcal P_{gg}\otimes\Pi]\big)-\frac{11}{12}+[\mathcal P_{gg}\otimes\Pi]=-\frac{\ell_\vee}{2}$, the same total as in the two-$\epsilon$ bookkeeping. The poles proportional to $L$ cancel within the JIMWLK sector in both bookkeepings. The remainder $-\frac{\ell_\vee}{2}\frac2\epsilon$ is a pole multiplied by $\ell_\vee=\log(1-k^+/\vee)=-k^+/\vee+O(k^{+2}/\vee^2)$, a power correction in the ratio of the soft and the valence longitudinal momenta. It is not cancelled, and we do not claim that it is: it is outside the eikonal approximation. $\ell_\vee=-\int_0^{k^+/\vee}d\xi/(1-\xi)$ is the part of the vertex correction (and, with the opposite sign of the pole, of the real splitting) in which the parent of the observed gluon is faster than $\vee$, i.e.\ a valence mode, and which is not enhanced at small momentum fraction (the enhanced part $2/\xi$ of the same integrand gives the $\log(\vee/k^+)$ that cancels in the pole coefficient). For this part the description of the valence modes by a static source is not valid, and its coefficient is not defined by the effective theory: for a smooth separation $f(p^+/\vee)$ of soft and valence modes ($f(0)=1$) $\ell_\vee$ becomes $-\kappa_fk^+/\vee$ with $\kappa_f=\int_0^\infty dv\,[1-f(v)]/v^2$ ($\kappa_f=1,\sqrt\pi,\pi/2$ for a sharp cut-off, $e^{-u^2}$ and $1/(1+u^2)$), while $\log(\vee/k^+)$ only shifts by a universal constant (Section~\ref{sec:fin_inst}). We therefore define the eikonal approximation as the leading term of the expansion in $k^+/\vee$ and $\Lambda/k^+$, with the logarithms kept, taken at the level of the longitudinal integrands at fixed transverse momenta and at fixed $\epsilon$, i.e.\ before the transverse integrals are done and the regulator is removed. Since the pole factor of Row~I.II is a product of a longitudinal and a transverse factor, $\big(\frac{11}{3}-2L+\ell_\vee\big)\big[-\frac2\epsilon+\log\frac{\v k^2}{\mu^2}\big]$, this expansion is uniform in the transverse loop momentum; at fixed $\epsilon$ it is equivalent to $\vee\to\infty$ before $\epsilon\to0$. A pole multiplying a power correction would have to be renormalised, at that power, by non-eikonal interactions of the valence sector that are not part of the Hamiltonian~\eqref{eq:Hdecomp}; at leading power it is absent. The same holds for the terms $O(\Lambda/k^+)$ on the target side.

The account answers the questions raised by the virtual $\frac{11}{3}$ terms. (i) The virtual term of the vertex correction is present, Eq.~\eqref{eq:dg_II}, and is cancelled, ultraviolet against ultraviolet, by Eq.~\eqref{eq:dg_V}; hence it does not appear in the final result. (ii) The pole of the term $\frac{\beta_0}{2}\delta(1-\xi)$ as it arises in Row~II.IV is collinear. What renormalises the coupling is the ultraviolet half of the normalisation of the dressed produced gluon, $Z_3$ in light-cone gauge, consistent with the Ward identity $Z_g=Z_3^{-1/2}$ ($g_0A_0=gA$, so that the coupling of the classical source to the field needs no counterterm of its own): after the cancellation of Rows~I.II and~I.V the only ultraviolet pole left is that of $Z_3$, $+\frac{\alpha_s\beta_0}{4\pi}\frac{2}{\epsilon_{\rm UV}}$, which is what $g_0^2\to g^2(\mu)$ removes. (iii) The infrared half of the same normalisation is the virtual $\frac{\beta_0}{2}\delta(1-\xi)$ (and $-2N_cL\,\delta(1-\xi)$) of the final-state splitting; real $\mathcal P_{gg+}$ and virtual term together are the full $\mathcal P_{gg}$, which is absorbed by the fragmentation counterterm. No infrared pole is left. In the single-$\epsilon$ bookkeeping the coupling counterterm cancels the collinear pole of the $\delta(1-\xi)$ term formally; this is legitimate only because the normalisation integral is scaleless.

\paragraph{Remark (pure Yang--Mills; identification of the $\beta_0$ terms).} The calculation and the result are for pure Yang--Mills theory, $N_f=0$, where $\beta_0/(2N_c)=\frac{11}{6}$ coincides with integrals of the splitting function. Which of the numbers $\frac{11}{6}$ are $\beta_0/(2N_c)$ can be decided by the BLM-type device of adding a quark loop, without computing the $N_f$ terms of the result, which we do not give. With the vertices normalised as here (the paper's triple-gluon matrix element is, exactly, minus the covariant vertex contracted with light-cone polarisation vectors, and the quark vertex is taken in the same normalisation), the angular-averaged squared vertices satisfy $\frac{n_fT_R\sum|V_{q\bar q}|^2}{\frac12N_c\sum|V_{gg}|^2}=\frac{n_fT_R}{N_c}\frac{\xi^2+(1-\xi)^2}{K_{gg}(\xi)}$, $K_{gg}=\xi(1-\xi)-2+\frac{1}{\xi(1-\xi)}=\mathcal P_{gg+}/(2N_c)$, for every $\xi$: a $q\bar q$ loop adds to the kernel $K_{gg}$ of the self-energy the kernel $\frac{n_fT_R}{N_c}[\xi^2+(1-\xi)^2]$ (with its $d_\perp$-dimensional spin trace), and the corresponding $q\bar q$ pair crosses the shock wave in Row~I.V, with colour factor $T_R[U(\v y)-U(\v x)]$ for a point-like pair. Hence the $\frac{11}{3}$ of Row~I.II becomes $\beta_0/N_c$, $\beta_0=\frac{11}{3}N_c-\frac43T_Rn_f$, the $-\frac{11}{6}$ of $N_{2b\,\rm UV}$ becomes $-\beta_0/(2N_c)$, the $\frac{11}{6}$ of the subtraction term in $\mathfrak K^{\rm V}_{2b}$ becomes $\beta_0/(2N_c)$, the $\frac{11}{12}$ of the normalisation, of the coupling counterterm and of the RC term become $\beta_0/(4N_c)$, and the $\frac{\beta_0}{2}$ of $\delta(1-\xi)$ becomes the full $\frac{\beta_0}{2}$: all of these are integrals of the complete one-loop gluon self-energy kernel. In contrast, the coefficient $-\frac{11}{6}$ of the large-distance tail, Eq.~\eqref{eq:dg_tailtot}, is the double-plus integral of the real $\mathcal P_{gg+}$ and stays $-\frac{11}{6}$ for any $N_f$ (the $\delta(1-\xi)$ parts of the DGLAP and RC terms cancel in the tail for any $\beta_0$). The constant of Row~I.II would become $\frac{67}{9}-\frac{\pi^2}{3}-\frac{20}{9}\frac{n_fT_R}{N_c}=\frac{2}{N_c}K_{\rm CMW}$, $K_{\rm CMW}=N_c\big(\frac{67}{18}-\frac{\pi^2}{6}\big)-\frac{10}{9}T_Rn_f$. There is no real $N_f$ term for a measured gluon at this order. The cancellation of the ultraviolet poles between Rows~I.II and I.V, and the pole balance of Table~\ref{tab:poles}, hold for any $\beta_0$.

\subsection{The collinear term in the \texorpdfstring{$\MS$}{MS-bar} scheme}\label{sec:dg_msbar}

\paragraph{The bare collinear term.} In the variables of Row~II.IV, $\v r=\v w-\v z$, $\v u=\v w'-\v w$ and $\v Y=\frac{\v y+\v y'}{2}$, with $\v y'-\v y=\xi\v u$, the Jacobian of $(\v z,\v w,\v w')\to(\v Y,\v r,\v u)$ is one in any dimension, and the $\v r$ integral gives $-\pi\Gamma(1-\frac\epsilon2)\big(\frac{\mu^2\v u^2e^{\gamma_E}}{4}\big)^{\epsilon/2}\big[\frac{\delta^{mm'}}{\epsilon}+\frac{u^mu^{m'}}{\v u^2}\big]$. If $\v u$ is continued to $d$ dimensions as well and the separation of the parents, $\v s=\v y'-\v y=\xi\v u$, is used as variable,
\begin{align}\label{eq:dg_resc}
\int d^du\,(\v u^2)^{\epsilon/2}\,e^{i\v k\cdot\v u}\,f(\xi\v u)=\xi^{-d}\xi^{-\epsilon}\int d^ds\,(\v s^2)^{\epsilon/2}\,e^{i\frac{\v k}{\xi}\cdot\v s}\,f(\v s)=\xi^{-2}\int d^ds\,(\v s^2)^{\epsilon/2}\,e^{i\frac{\v k}{\xi}\cdot\v s}\,f(\v s)\,.
\end{align}
The bare collinear term is therefore $\xi^{-2}\big[\frac2\epsilon+\mathcal B(\v y'-\v y)\big]_d$ times the LO structure of the parent, without $\log\xi^2$ in the logarithm; a $\log\xi^2$ arises if only $\v r$ is continued to $d$ dimensions, and the two sides of Eq.~\eqref{eq:dg_resc} then differ by $\xi^{-\epsilon}$.

\paragraph{Definition of the fragmentation function and the counterterm.} What survives in the subtracted result depends on the power of $\xi$ that multiplies the counterterm, i.e.\ on the definition of the fragmentation function. We define the fragmentation function as an $\MS$ number density in the momentum fraction (operator definition, momentum sum rule $\sum_h\int dz\,z\,D_h(z)=1$ in $d$ dimensions). Collinear kinematics $\v P=z\v p$ in all $d$ transverse components then gives
\begin{align}
\frac{dN_h}{dP^+d^dP}=\int\frac{dz}{z^{1+d}}\,D(z)\,\frac{dN_g}{dk^+d^dk}\bigg|_{k=P/z}\,,
\end{align}
and the $\MS$ renormalisation $D_0=\big[\delta(1-\xi)+\frac{\alpha_s}{2\pi}\frac2\epsilon\mathcal P_{gg}(\xi)+\dots\big]\otimes D(\mu)$ (pure pole, Mellin convolution) gives the partonic counterterm $+\mathcal C\frac2\epsilon\int d\xi\,\mathcal P_{gg}(\xi)\,\xi^{-d}\,\widehat L_d(\v k/\xi)$, with $\xi^{-d}$ rather than $\xi^{-2}$. This is the counterterm used throughout. It is the standard $\MS$ convention: for a cross section differential in all components of the momentum of the observed hadron the $\MS$ collinear counterterm carries $dz/z^{2-2\epsilon}$ in $4-2\epsilon$ dimensions, i.e.\ $\xi^{-d}$ here, Eq.~(6.17) of~\cite{Catani:1996vz}. With it, and with the same pure-pole renormalisation, the one-hadron-inclusive gluon coefficient functions of $e^+e^-\to gX$~\cite{Altarelli:1979kv,Kretzer:2000yf} and, for $g\to g$ in pure Yang--Mills, of $H\to gX$~\cite{Almasy:2011eq} are reproduced exactly (we computed both, with the $d$-dimensional real-emission matrix elements, phase space and counterterm, in exact arithmetic). With $\xi^{-2}$ in place of $\xi^{-d}$ both come out wrong by $-\frac{\alpha_s}{2\pi}P(z)\log z^2$ times the Born term, the structure of $\Delta$ in Eq.~\eqref{eq:dg_Delta} below. These checks are angle-integrated; for the transverse-momentum Jacobian they rely on the universality of the fragmentation function and on $d$-dimensional collinear kinematics, which is the derivation of~\cite{Catani:1996vz}. With $\xi^{-d}=\xi^{-2}(1+\epsilon\log\xi)$ (CT: counterterm),
\begin{align}\label{eq:dg_two}
\text{bare}-\text{CT}(\xi^{-d})&\to\xi^{-2}\!\int e^{-i\frac{k}{\xi}\cdot(y-y')}\,\log\frac{(y'-y)^2\mu^2}{c_0^2\,\xi^2}\,\mathcal O(y,y')\,,\notag\\
\text{bare}-\text{CT}(\xi^{-2})&\to\xi^{-2}\!\int e^{-i\frac{k}{\xi}\cdot(y-y')}\,\log\frac{(y'-y)^2\mu^2}{c_0^2}\,\mathcal O(y,y')\,.
\end{align}
The first line is the collinear term for the fragmentation function defined above, and is the form used in this paper: the argument of its logarithm is the transverse separation $(\v y'-\v y)/\xi$ of the observed gluon in the amplitude and in the conjugate amplitude, the variable conjugate to $\v k$. The second line is the form of Chirilli, Xiao and Yuan (CXY)~\cite{Chirilli:2011km,Chirilli:2012jd} and of the $\MS$ form of~\cite{Altinoluk:2014eka}, with the separation of the parent. Both statements were verified numerically in $d=2\pm0.002$ with $d$-dimensional Fourier--Bessel transforms of a test function, to $O(\epsilon^2)$. Hence the $\log\xi^2$ cancels in the bare result, but not in the subtracted one, since with our definition the counterterm carries the $d$-dimensional Jacobian. CXY obtain the collinear pole of the final-state splitting from a $d$-dimensional integral over the momentum of the parent at fixed observed momentum, which yields exactly $\xi^{-2}$ times the parent spectrum, and absorb it by a pure-pole renormalisation of $D_{h/g}$ in a hadronic formula written with $dz/z^2$~\cite{Chirilli:2012jd}: their counterterm carries $\xi^{-2}$. The appendix of~\cite{Altinoluk:2014eka} writes the hadronic formula with the $d$-dimensional Jacobian $\zeta^{-(D-2)}$, but defines the renormalised fragmentation function with the measure $d\xi/(1-\xi)^{5-D}$, which is not the pure-pole renormalisation of the number density and amounts to the same choice as CXY. In the cut-off scheme of the main text of~\cite{Altinoluk:2014eka} our form is obtained with $\mu^2_{\MS}=c_0^2\mu^2_{\rm cut}\approx1.26\,\mu^2_{\rm cut}$.

\paragraph{Remark: the CXY convention.} With $\mathcal G_0(\xi)\equiv\xi^{-2}\int_{x,y,x',y'}e^{-i\v k\cdot(\v y-\v y')/\xi}\,\mathcal O(y,y')=\Pi(\xi)$, the two forms of the collinear term differ by
\begin{align}\label{eq:dg_Delta}
\Delta\equiv\frac{d^3N_{\rm DGLAP}}{d^3k}-\frac{d^3N^{\rm CXY}_{\rm DGLAP}}{d^3k}=+\,\mathcal C\int_0^1d\xi\,\mathcal P_{gg}(\xi)\,\log\xi^2\;\mathcal G_0(\xi)\,.
\end{align}
The $\delta(1-\xi)$ part of $\Delta$ vanishes, $\Delta$ is integrable at $\xi\to0$ ($\mathcal G_0(\xi)\to0$), and it does not depend on $\mu$. Equivalently $D^{\rm CXY}=\big[\delta(1-\xi)+\frac{\alpha_s}{2\pi}\mathcal P_{gg}(\xi)\log\xi^2\big]\otimes D$. The results of CXY and of~\cite{Altinoluk:2014eka} are meant to be used with standard $\MS$ fragmentation functions, which obey the momentum sum rule and for which the $\MS$ coefficient is the first line of Eq.~\eqref{eq:dg_two}; their form of the collinear term is the coefficient appropriate for the functions $D^{\rm CXY}$, i.e.\ it misses $\Delta$ at $O(\alpha_s)$ when used with standard fragmentation functions. Since $\int_0^1d\xi\,\xi\,\mathcal P_{gg}(\xi)\log\xi^2=4N_c\big(\frac{65}{144}-\zeta_2\big)=\frac{N_c(65-24\pi^2)}{36}\neq0$ ($-14.3225$ for $N_c=3$; with quarks the combination $\int_0^1d\xi\,\xi\log\xi^2[\mathcal P_{gg}+2N_f\mathcal P_{qg}]=\frac{N_c(65-24\pi^2)}{36}-\frac{11N_f}{36}$ would enter), the functions $D^{\rm CXY}$ do not satisfy the momentum sum rule at $O(\alpha_s)$; equivalently, the energy sums of the $e^+e^-$ and $H\to gX$ coefficient functions computed with $\xi^{-2}$ are not those of the total rates. $\Delta$ is given only so that the CXY convention can be recovered; all collinear terms of this paper refer to the first line of Eq.~\eqref{eq:dg_two}. A comparison with the NLO coefficient of hadroproduction, $pp\to gX$, and a collinear (large-$\v k$) limit of the hybrid formula of~\cite{Chirilli:2012jd} have not been carried out; the bare final-state collinear terms of~\cite{Chirilli:2012jd} and of this paper are the same $d$-dimensional integral, so that such a comparison would re-test only the counterterm, which is fixed above.

\subsection{The DGLAP and running-coupling terms}\label{sec:dg_final}

We give the DGLAP and running-coupling terms in two exactly equivalent forms: in coordinate space, Eqs.~\eqref{eq:DGLAP}--\eqref{eq:LORC}, and in a local momentum-space form, Eqs.~\eqref{eq:dg_local}--\eqref{eq:dg_R}; their relation is Eq.~\eqref{eq:dg_equiv}. Both refer to the $\MS$ number-density fragmentation function of Section~\ref{sec:dg_msbar}.

\paragraph{Coordinate-space form.} With one scale $\mu$ and the phase $e^{-i\v k\cdot(\v y-\v y')/\xi}$, define
\begin{align}\label{eq:dg_F}
\mathcal F(\xi)\equiv\frac{1}{\xi^2}\int_{x,y,x',y'}e^{-ik\cdot(y-y')\frac1\xi}\,\log\frac{(y'-y)^2\mu^2}{c_0^2\,\xi^2}\;\mathcal O(y,y')\,,\qquad
\mathcal F(0)\equiv\lim_{\xi\to0}\mathcal F(\xi)=-\frac{4\pi}{\v k^2}\,\mathcal X\,,
\end{align}
where the integral over $\v y-\v y'$ is done before the integral over $\xi$ ($\mathcal F(0)$ is the same in the CXY form). The DGLAP term is
\begin{align}\label{eq:DGLAP}
\frac{d^3N_{\rm DGLAP}}{d^3k}&=-\mathcal C\int_0^1d\xi\,\widetilde{\mathcal P}_{gg}(\xi)\,\mathcal F(\xi)\,,\notag\\
\widetilde{\mathcal P}_{gg}(\xi)&=2N_c\Big[\xi(1-\xi)+\frac{\xi}{(1-\xi)_+}+\Big[\frac{1-\xi}{\xi}\Big]_+\Big]+\frac{\beta_0}{2}\,\delta(1-\xi)\,,\notag\\
\frac{d^3N_{\rm DGLAP}}{d^3k}&=-2N_c\,\mathcal C\int_0^1d\xi\Big\{\xi(1-\xi)\,\mathcal F(\xi)+\frac{\xi\mathcal F(\xi)-\mathcal F(1)}{1-\xi}+\frac{1-\xi}{\xi}\big[\mathcal F(\xi)-\mathcal F(0)\big]\Big\}\notag\\&\quad-\mathcal C\,\frac{\beta_0}{2}\,\mathcal F(1)\,,
\end{align}
with the prescriptions~\eqref{eq:plusxi}. $\widetilde{\mathcal P}_{gg}$ is the $\MS$ splitting function $\mathcal P_{gg}$ with, in addition, the subtraction at $\xi\to0$, whose subtracted term belongs to the BFKL term of the row. With this reading the left-over of that subtraction,
\begin{align}\label{eq:remnant}
\frac{d^3N_{\rm rem}}{d^3k}=-2N_c\,\mathcal C\,\frac{4\pi}{\v k^2}\,\mathcal X=-\frac{1}{(2\pi)^3}\,\frac{g^4N_c}{2\pi^3}\,\frac{1}{\v k^2\,k^+}\,\mathcal X\,,
\end{align}
the term $-\frac{1}{k^+}\times$(BFKL structure), is counted among the finite terms of Row~II.IV (Section~\ref{sec:finite}). It arises because $\int_{k^+/\vee}^1d\xi\frac{1-\xi}{\xi}=\log\frac{\vee}{k^+}-1+\frac{k^+}{\vee}$. The sum of Eqs.~\eqref{eq:DGLAP} and~\eqref{eq:remnant} can be written with the plus prescription at $\xi\to0$ acting on $(1-\xi)\mathcal F(\xi)$,
\begin{align}\label{eq:DGLAPdoubleplus}
\frac{d^3N_{\rm DGLAP}}{d^3k}+\frac{d^3N_{\rm rem}}{d^3k}&=-\mathcal C\int_0^1d\xi\,\Big\{\mathcal P_{gg}(\xi)\,\mathcal F(\xi)-\frac{2N_c}{\xi}\,\mathcal F(0)\Big\}\notag\\
&=-2N_c\,\mathcal C\int_0^1\frac{d\xi}{\xi_+(1-\xi)_+}\,N(\xi)\,\mathcal F(\xi)-\mathcal C\,\frac{\beta_0}{2}\,\mathcal F(1)\,,
\end{align}
where $\mathcal P_{gg}$ has no subtraction at $\xi=0$ and the last form is the ``double plus'' prescription of~\cite{Kovner:2023vsy} (Appendix~\ref{app:distributions}); the two subtractions are there for the same reason as in~\cite{Kovner:2023vsy}: both poles are part of the leading-logarithmic evolution, here BFKL at $\xi\to0$ and JIMWLK at $\xi\to1$. Eqs.~\eqref{eq:DGLAP}--\eqref{eq:remnant} and Eq.~\eqref{eq:DGLAPdoubleplus} are the same quantity; only one of the two forms is to be used. The dependence on the scale is $\partial N_{\rm DGLAP}/\partial\log\mu^2=-\mathcal C\int_0^1d\xi\,\mathcal P_{gg}(\xi)\,\Pi(\xi)$, the standard DGLAP evolution of $D_g(\mu)$ including its $\frac{\beta_0}{2}\delta(1-\xi)$; the subtraction at $\xi=0$ does not affect it, since $\Pi(0)=0$.

The running-coupling term is the $\delta(1-\xi)$ part of Eq.~\eqref{eq:DGLAP} with the opposite sign,
\begin{align}\label{eq:RC}
\frac{d^3N_{\rm RC}}{d^3k}=+\,\mathcal C\,\frac{\beta_0}{2}\,\mathcal F(1)=+\frac{\alpha_s\beta_0}{4\pi}\,\mathcal N_0\int_{x,y,x',y'}e^{-ik\cdot(y-y')}\,\log\frac{(y'-y)^2\mu^2}{c_0^2}\;\mathcal O(y,y')\,.
\end{align}
It is added and subtracted: $N_{\rm DGLAP}+N_{\rm RC}$ does not depend on $\beta_0$, since the real emission produces no $\delta(1-\xi)$ and the virtual $\frac{11}{3}$ terms cancel between Rows~I.II and~I.V. Its content is the dependence on the scale: $\partial N_{\rm RC}/\partial\log\mu^2=+\frac{\alpha_s\beta_0}{4\pi}\times$LO compensates the running of $g^2(\mu)$ in the LO term, and the $\delta(1-\xi)$ part of $N_{\rm DGLAP}$ the corresponding part of the evolution of $D_g(\mu)$. Together with the LO term,
\begin{align}\label{eq:LORC}
\frac{d^3N_{\rm LO}}{d^3k}+\frac{d^3N_{\rm RC}}{d^3k}&=\frac{1}{(2\pi)^3}\,\frac{g^2(\mu)}{2\pi^2}\,\frac{1}{k^+}\notag\\&\times\int_{x,y,x',y'}e^{-ik\cdot(y-y')}\Big[1+\frac{\alpha_s(\mu)\beta_0}{4\pi}\log\frac{(y'-y)^2\mu^2}{c_0^2}\Big]\mathcal O(y,y')\,,
\end{align}
and to this order the bracket is $\alpha_s(c_0^2/\v r^2)/\alpha_s(\mu^2)$, $\v r=\v y'-\v y$. Since $N_{\rm RC}$ is minus the $\delta(1-\xi)$ term of $N_{\rm DGLAP}$, the scale $c_0^2/\v r^2$ is the choice $\mu=c_0/r$ in the $\delta(1-\xi)$ part of the collinear logarithm, the choice made by CXY for the factorisation scale~\cite{Chirilli:2012jd}; Eq.~\eqref{eq:LORC} holds only together with the $\delta(1-\xi)$ term kept in $\widetilde{\mathcal P}_{gg}$ at the same $\mu$.

\paragraph{Separate renormalisation and factorisation scales.} The add-and-subtract of the $\delta(1-\xi)$ term identifies the renormalisation scale with the factorisation scale; the statement that the residual dependence on $\mu$ comes only from the $\delta(1-\xi)$ term holds only for $\mu_R=\mu_F$. With $g_0^2=\mu_R^\epsilon g^2(\mu_R)\big[1-\frac{\alpha_s\beta_0}{4\pi}\frac2\epsilon\big]$ and $D_0=\big[\delta(1-\xi)+\frac{\alpha_s}{2\pi}\frac2\epsilon\mathcal P_{gg}\big]\otimes D(\mu_F)$, the two halves of the scaleless normalisation~\eqref{eq:dg_norm} are cancelled separately, the ultraviolet half by the coupling counterterm at $\mu_R$ and the infrared half by the $\delta(1-\xi)$ part of the fragmentation counterterm at $\mu_F$, and
\begin{align}\label{eq:dg_muRmuF}
\frac{d^3N_{\rm LO}}{d^3k}+\frac{d^3N_{\rm RC}}{d^3k}&=\frac{1}{(2\pi)^3}\,\frac{g^2(\mu_R)}{2\pi^2}\,\frac{1}{k^+}\notag\\&\quad\times\int_{x,y,x',y'}e^{-ik\cdot(y-y')}\Big[1+\frac{\alpha_s\beta_0}{4\pi}\log\frac{(y'-y)^2\mu_R^2}{c_0^2}\Big]\mathcal O(y,y')\,,\notag\\
\frac{d^3N_{\rm DGLAP}}{d^3k}&=-\mathcal C\int_0^1d\xi\,\widetilde{\mathcal P}_{gg}(\xi)\,\mathcal F(\xi)\Big|_{\mu\to\mu_F}\,,
\end{align}
including the $\frac{\beta_0}{2}\delta(1-\xi)$ term at $\mu_F$. Their $\beta_0$ part is the running-coupling term at $\mu_R$ minus the $\delta(1-\xi)$ term at $\mu_F$, $N_{\rm RC}(\mu_R)-\mathcal C\frac{\beta_0}{2}\mathcal F(1)|_{\mu_F}$, which is $+\frac{\alpha_s\beta_0}{4\pi}\log\frac{\mu_R^2}{\mu_F^2}\times$LO, and the result is independent of $\mu_R$ and $\mu_F$ to this order. In the local form, $N^{\rm loc}_{\rm RC}=-\frac{\alpha_s\beta_0}{4\pi}\log\frac{\v k^2}{\mu_R^2}\,\mathcal N_0\widehat L(k)$ and $N^{\rm loc}_{\rm DGLAP}=\mathcal C\log\frac{\v k^2}{\mu_F^2}\int_0^1d\xi\,\mathcal P_{gg}(\xi)\,\Pi(\xi)$, while $\mathcal R$ and $N_{\mathcal R}$ are unchanged.

In the classification of~\cite{Kovner:2023vsy}, where of the two $\beta_0$'s of the NLO JIMWLK kernel one is the ultraviolet running of the coupling and the other a DGLAP-type resolution logarithm, both are here tied to the produced gluon: the ultraviolet one is the $Z_3$ pole of the scaleless normalisation~\eqref{eq:dg_norm}, removed by the coupling counterterm and not visible in the finite result; the infrared one is the $\frac{\beta_0}{2}\delta(1-\xi)$ of $\mathcal P_{gg}$, the only $\beta_0$ that appears explicitly. They are equal in magnitude because the normalisation integral is scaleless. This agrees in structure with Kovchegov and Weigert~\cite{Kovchegov:2007vf} and Horowitz and Kovchegov~\cite{Horowitz:2010yg}, in whose calculation (quark--quark scattering, BLM resummation of quark bubbles) the observable includes collinear pairs whose invariant mass, i.e.\ the virtuality of the produced-gluon line, is below a fixed cut-off $\Lambda_{\rm coll}$, and the coupling of the emission vertex of the produced gluon is $\alpha_s(\Lambda^2_{\rm coll}e^{-5/3})$ (Eqs.~(3.10)--(3.16) of~\cite{Horowitz:2010yg}); $\Lambda_{\rm coll}$ does not depend on the exchanged momenta and plays the role of the factorisation scale of the fragmentation function~\cite{Kovchegov:2007vf}, while the remaining couplings run with $\v q^2$, $(\v k-\v q)^2$ and a complex scale $Q(\v q,\v k-\v q)$, Eqs.~(3.31), (3.32) of~\cite{Horowitz:2010yg}. (Equation numbers of~\cite{Horowitz:2010yg} refer to the arXiv version.)

\paragraph{Local (momentum-space) form.} With $\log\frac{\v r^2\mu^2}{c_0^2}=-\frac1\pi\int d^2l\,\frac{e^{-i\v l\cdot\v r}-\theta(\mu-|\v l|)}{\v l^2}$ one has exactly
\begin{align}\label{eq:dg_loc}
\mathcal F(\xi)=-\log\frac{\v k^2}{\mu^2}\,\Pi(\xi)+\frac{1}{\xi^2}\,\mathcal R\Big(\frac{\v k}{\xi}\Big)\,,\qquad
\mathcal R(\v q)\equiv-\frac1\pi\int d^2l\,\frac{\widehat L(\v q+\v l)-\theta(|\v q|-|\v l|)\,\widehat L(\v q)}{\v l^2}\,,
\end{align}
with $\lim_{\xi\to0}\xi^{-2}\mathcal R(\v k/\xi)=\mathcal F(0)$; $\mathcal R$ does not depend on $\mu$. Inserted into Eqs.~\eqref{eq:DGLAP}--\eqref{eq:RC}, it splits each term into a local part and a non-local, $\mu$-independent part. The local DGLAP and running-coupling terms are
\begin{align}\label{eq:dg_local}
\frac{d^3N^{\rm loc}_{\rm DGLAP}}{d^3k}&\equiv\mathcal C\,\log\frac{\v k^2}{\mu^2}\int_0^1d\xi\,\mathcal P_{gg}(\xi)\,\Pi(\xi)\,,\qquad
\frac{d^3N^{\rm loc}_{\rm RC}}{d^3k}\equiv-\frac{\alpha_s\beta_0}{4\pi}\,\log\frac{\v k^2}{\mu^2}\,\mathcal N_0\,\widehat L(k)\,,\notag\\
\frac{d^3N_{\rm LO}}{d^3k}+\frac{d^3N^{\rm loc}_{\rm RC}}{d^3k}&=\mathcal N_0\,\widehat L(k)\Big[1-\frac{\alpha_s\beta_0}{4\pi}\log\frac{\v k^2}{\mu^2}\Big]\,,
\end{align}
and the coordinate-space terms are
\begin{align}\label{eq:dg_locsplit}
\frac{d^3N_{\rm DGLAP}}{d^3k}+\frac{d^3N_{\rm rem}}{d^3k}&=\frac{d^3N^{\rm loc}_{\rm DGLAP}}{d^3k}-\mathcal C\int_0^1d\xi\Big\{\mathcal P_{gg}(\xi)\,\frac{\mathcal R(\v k/\xi)}{\xi^2}-\frac{2N_c}{\xi}\,\mathcal F(0)\Big\}\,,\notag\\
\frac{d^3N_{\rm RC}}{d^3k}&=\frac{d^3N^{\rm loc}_{\rm RC}}{d^3k}+\mathcal C\,\frac{\beta_0}{2}\,\mathcal R(\v k)\,.
\end{align}
The local terms carry the whole dependence on $\mu$ and need no subtraction at $\xi=0$ ($\Pi(\xi)\to0$); to this order $N_{\rm LO}+N^{\rm loc}_{\rm RC}$ is the LO spectrum with the coupling at the measured momentum, $\alpha_s(\v k^2)/\alpha_s(\mu^2)$. The non-local remainders of the two lines of Eq.~\eqref{eq:dg_locsplit}, whose $\delta(1-\xi)$ parts $\mp\mathcal C\frac{\beta_0}{2}\mathcal R(\v k)$ cancel, add up to
\begin{align}\label{eq:dg_R}
\frac{d^3N_{\mathcal R}}{d^3k}=-\mathcal C\int_0^1d\xi\Big\{\mathcal P_{gg+}(\xi)\,\frac{\mathcal R(\v k/\xi)}{\xi^2}-\frac{2N_c}{\xi}\,\mathcal F(0)\Big\}\,,
\end{align}
which belongs to the finite terms of Row~II.IV. It contains the BFKL ($\xi\to0$) subtraction and the region of a soft parent discussed next. The two forms are related by
\begin{align}\label{eq:dg_equiv}
\frac{d^3N_{\rm DGLAP}}{d^3k}+\frac{d^3N_{\rm rem}}{d^3k}+\frac{d^3N_{\rm RC}}{d^3k}=\frac{d^3N^{\rm loc}_{\rm DGLAP}}{d^3k}+\frac{d^3N^{\rm loc}_{\rm RC}}{d^3k}+\frac{d^3N_{\mathcal R}}{d^3k}\,,
\end{align}
an identity; either side may be used, with the same finite terms of the other rows.

\paragraph{The two forms at large distance.} The forms differ in where the large-distance (soft-parent) region of Row~II.IV sits (Section~\ref{sec:dg_tail}). In the coordinate-space form the DGLAP and running-coupling terms (and the remnant) are each logarithmically divergent at large position of the parent gluon. Their sum with the finite terms of the second part of Row~II.I and of Row~III.III is not: the tail cancels term by term between the three rows (Table~\ref{tab:tails}), which therefore have to be integrated together over the position of the gluon pair. In the local form $N^{\rm loc}_{\rm DGLAP}$ and $N^{\rm loc}_{\rm RC}$ have no large-distance tail; the tail of the coordinate-space DGLAP, running-coupling and remnant terms together, $-\frac{11}{6}\,\mathcal U$ in the units of Table~\ref{tab:tails}, is carried by the remainder $N_{\mathcal R}$, which is then integrated together with the second part of Row~II.I and with Row~III.III. In both forms the $\cos2\varphi$ part of the tail of Row~II.IV is in the tensor term of its finite part.

\subsection{The large-distance tail of Row II.IV}\label{sec:dg_tail}

At large parent position $\v Y$ the LO structure tends to $G_\infty(\v Y)/\v Y^2$, $G_\infty=\int_{x,x'}[U(\v Y)-U(\v x')]^{ce'}[U(\v Y)-U(\v x)]^{ce}\rho^{e'}(\v x')\rho^e(\v x)$, independently of the separation $\v s=\v y'-\v y$. The logarithm then survives the Fourier transform, $\int d^2s\,e^{-i\v q\cdot\v s}\log\v s^2=-4\pi/\v q^2$, and $\mathcal F(\xi)\to-\frac{4\pi}{\v k^2}G_\infty/\v Y^2$ for every $\xi$. In units of $\mathcal C\frac{4\pi}{\v k^2}\int\frac{d^2Y}{\v Y^2}G_\infty$ the tails of the coordinate-space terms are: DGLAP $+\frac{N_c}{6}$, RC $-\frac{11N_c}{6}$, remnant $-2N_c$, and the traceless tensor term of the finite part of the row, Eq.~\eqref{eq:H-II.IVrem}, $+\frac{N_c}{3}\cos2\varphi$, $\varphi$ the angle between $\v k$ and $\v Y$; in total
\begin{align}\label{eq:dg_tailtot}
2N_c\,\mathcal C\,\frac{4\pi}{\v k^2}\int\frac{d^2Y}{\v Y^2}\Big[-\frac{11}{6}+\frac16\cos2\varphi\Big]\,G_\infty(\v Y)\,.
\end{align}
The coefficient $-\frac{11}{6}=\int_0^1\frac{d\xi}{\xi_+(1-\xi)_+}N(\xi)$ is the double-plus integral of the real splitting function. The RC term and the $\delta(1-\xi)$ part of the DGLAP term cancel each other also in the tail, so the tail is not related to the running coupling, nor to the pole ($\mathcal G_0$ has no tail for $\v k\neq0$). $G_\infty\neq0$ also for a colour-neutral projectile and in the colour-singlet projection.

\paragraph{Origin.} In momentum space the tail is the region $\v k/\xi+\v l\to0$ of $\mathcal R$: a parent much softer than the measured gluon, $|\v q|\ll|\v k|$, which splits into two hard gluons. Since the LO spectrum behaves as $\widehat L(\v q)\propto1/\v q^2$ at small $\v q$ (Weizs\"acker--Williams field of the colour charge $\int[U_\infty-U(\v x)]\rho(\v x)$, non-zero also for a neutral projectile), $\int d^2q\,\widehat L(\v q)/\v k^2$ diverges logarithmically. In a model with an infrared cut-off $m$ the coordinate-space form grows like $-\frac{2}{\v k^2}\log\frac1m$, while the local form $-\log\frac{\v k^2}{\mu^2}\widehat L(\v k)$ is independent of $m$. The DGLAP and RC terms as written in coordinate space are thus separately logarithmically divergent at large distance; the local terms~\eqref{eq:dg_local} are not, and the divergence resides in $N_{\mathcal R}$.

\paragraph{Cancellation.} The two-gluon final state of Group~II and~III contains the amplitude $\mathcal A_B-\mathcal A_{CA}$ (Eq.~\eqref{eq:A2gamp}), where $\mathcal A_B$ is the two-gluon component of the wave function scattered on the target, through the single-$\rho$ part of $\mf B_2$, and $\mathcal A_{CA}$ the LO gluon split by $\bar{\mf C}_1$ after the shock wave. $|\mathcal A_B|^2$ is the second part of Row~II.I, $|\mathcal A_{CA}|^2$ is Row~II.IV, and $-\mathcal A^\dagger_{CA}\mathcal A_B$ with its complex conjugate is Row~III.III. In LCPT, with $E_1$ the energy of the one-gluon (parent) state and $E_2$ that of the pair, the two amplitudes are $b=\frac{V_{\rho g}V_{ggg}}{E_1E_2}$ and $c\,A=\frac{V_{\rho g}V_{ggg}}{E_1(E_2-E_1)}$, so that
\begin{align}\label{eq:dg_bcA}
b-c\,A=-\frac{V_{\rho g}V_{ggg}}{E_2(E_2-E_1)}\,:
\end{align}
the parent pole $1/E_1\propto q^+/\v q^2$ cancels. With the colour factors $[U(\v w)U(\v z)-U(\v x)]$ of $\mathcal A_B$ (in the reduced form) and $[U(\v y)-U(\v x)]$ of $\mathcal A_{CA}$, the amplitude is $[U(\v w)U(\v z)-U(\v y)]\,b+[U(\v y)-U(\v x)]\,(b-cA)$. We verified this numerically with the explicit momentum-space forms of $\mf B_2$ and of $\mf C_1\mf A^{(1)}$, Eqs.~\eqref{eq:B2mom} and~\eqref{eq:C1mom}: at small parent momentum the two amplitudes grow like $1/|\v q|$ while their difference tends to a constant, and the coefficients of $1/\v q^2$ in $|\mathcal A_B|^2$, $|\mathcal A_{CA}|^2$ and $2\,{\rm Re}\,\mathcal A_{CA}^\dagger\mathcal A_B$ are equal, i.e.\ they enter with the weights $1,1,-2$. Consequently, the large-distance (soft-parent) tail of Row~II.IV is cancelled by the large-distance region of the second part of Row~II.I and of Row~III.III, in which both final-state gluons are far from the projectile together ($\v w,\v w',\v z\to\infty$ at fixed $\v w-\v z$). The cancellation is pointwise in the momentum fraction and in the relative momentum of the pair, and hence holds separately in the BFKL, JIMWLK and finite sectors, including the $\cos2\varphi$ part, provided the three rows are split with the same end-point subtractions. It requires $U(\v w)U(\v z)f=fU(\v y)$ for the pair at large distance, i.e.\ a Wilson line that varies slowly there, as for a target of finite size. For a colour-neutral projectile no tail survives even without this condition, since the remaining term $[U(\v w)U(\v z)-U(\v y)]\,b$ multiplies $\int_x\Kv(\v x-\v Y)\rho(\v x)$, which falls like $1/\v Y^2$. The cancellation does not require a change of the split into DGLAP and RC terms, but the coordinate-space terms~\eqref{eq:DGLAP}--\eqref{eq:RC} are separately divergent at large distance: with the local form~\eqref{eq:dg_local} DGLAP and RC terms are finite, and the remainder~\eqref{eq:dg_R} has to be integrated together with the finite terms of the second part of Row~II.I and of Row~III.III, under one integral over the position of the pair. 
\paragraph{Term-by-term check.} We have verified the cancellation on the finite terms of the three rows. In the joint limit the three rows have the same tail density in $t=p^+/k^+$,
\begin{align}\label{eq:dg_taildens}
\rho(t)=\frac{(t^2+t+1)^2+t^2\cos2\varphi}{t(1+t)^4}\,,
\end{align}
with colour factors $+N_c\Gamma$, $+N_c\Gamma$ and $-N_c\Gamma$ (the last twice, with its complex conjugate), so that the sum vanishes at every $t$. Integrated, in units of $\mathcal U\equiv2N_c\,\mathcal C\frac{4\pi}{\v k^2}\int\frac{d^2Y}{\v Y^2}G_\infty$, the tails of the rows as written are given in Table~\ref{tab:tails}.
\begin{table}[h]\centering
\begin{tabular}{lc}\hline
row & tail$/\,\mathcal U$\\\hline
II.IV: DGLAP $\frac1{12}$, RC $-\frac{11}{12}$, remnant $-1$, tensor $\frac16\cos2\varphi$ & $-\frac{11}{6}+\frac16\cos2\varphi$\\
II.I, second part & $-\frac{11}{6}+\frac16\cos2\varphi$\\
III.III, second part, with c.c. & $+\frac{11}{3}-\frac13\cos2\varphi$\\
II.I and III.III, first parts & $0$\\\hline
sum & $0$\\\hline
\end{tabular}
\caption{Large-distance (soft-parent) tails of the finite terms in the joint limit $\v w,\v w',\v z\to\infty$ at fixed $\v w-\v z$.}\label{tab:tails}
\end{table}
In the first parts the measured gluon on one side is tied to the projectile by the Fourier transform, and there is no tail. In the BFKL and JIMWLK sectors each row carries $(1+1-2)\times$ the same logarithm; the piece moved between sectors by the recoil rewriting of Row~III.III is the companion of its logarithmic term and together with it reproduces the unsubtracted term. The rule $\int d^2z/\v z^2=0$ plays no role here, since the poles it removes multiply $\delta^2(\v k)$. The colour factors were checked with random $SU(2)$ and $SU(3)$ Wilson lines: for Wilson lines that are uniform at large distance the cancellation is exact (to $10^{-14}$), and for a colour-neutral projectile only the terms bilinear in $U(\v x)$, $U(\v x')$ survive, with factors $+N_c,+N_c,-N_c$ independent of the far Wilson lines, so that the tail cancels for any target. For a coloured projectile and a target that is not uniform at large distance a remainder proportional to the total charge survives; it is the non-singlet term of the JIMWLK formula, not a defect of the finite terms. The region in which only part of the pair is far does not contribute at this order (Section~\ref{sec:fin_integrable}, condition (C10)). Not verified are the interchange of the $p^+$ integral with the joint limit in Row~III.III (checked numerically for Row~II.I only) and a full transverse integration with Wilson lines.
